% !TeX program = pdflatex
\documentclass[12pt,a4paper]{article}

% Overleaf-friendly preamble. Compiler: pdfLaTeX.
\usepackage[T2A]{fontenc}
\usepackage[utf8]{inputenc}
\usepackage[english,russian]{babel}
\usepackage{cmap}

\usepackage{geometry}
\geometry{left=22mm,right=22mm,top=20mm,bottom=22mm}

\usepackage{amsmath,amssymb,mathtools,bm}
\usepackage{booktabs,longtable,array,tabularx}
\usepackage{enumitem}
\usepackage{xcolor}
\usepackage[unicode]{hyperref}
\usepackage{tcolorbox}
\tcbuselibrary{breakable}

\hypersetup{
  colorlinks=true,
  linkcolor=blue!45!black,
  urlcolor=blue!45!black,
  citecolor=blue!45!black
}

\setlist[itemize]{leftmargin=*, itemsep=2pt, topsep=4pt}
\setlist[enumerate]{leftmargin=*, itemsep=2pt, topsep=4pt}
\sloppy

\newcommand{\code}[1]{\texttt{#1}}
\newcommand{\R}{\mathbb{R}}
\newcommand{\E}{\mathbb{E}}
\newcommand{\Prob}{\mathbb{P}}
\newcommand{\Manifold}{\mathcal{M}}
\newcommand{\Hyp}{\mathcal{H}}
\newcommand{\dd}{\,\mathrm{d}}
\DeclareMathOperator{\clip}{clip}
\DeclareMathOperator{\diag}{diag}
\DeclareMathOperator{\covered}{covered}
\DeclareMathOperator{\conflict}{conflict}
\DeclareMathOperator{\Score}{Score}
\DeclareMathOperator{\norm}{norm}
\DeclareMathOperator{\Risk}{Risk}
\DeclareMathOperator{\Cov}{Coverage}

\newtcolorbox{mainbox}{
  breakable, colback=blue!3, colframe=blue!35!black,
  arc=1mm, boxrule=0.5pt, left=2mm, right=2mm, top=1mm, bottom=1mm
}

\newtcolorbox{formulaBox}{
  breakable, colback=gray!5, colframe=gray!55,
  arc=1mm, boxrule=0.5pt, left=2mm, right=2mm, top=1mm, bottom=1mm
}

\title{\textbf{Математическая модель системы\\«Фабрика гипотез»}}
\author{Диагностика потерь металлов и ранжирование гипотез}
\date{}

\begin{document}

\maketitle

\begin{mainbox}
\textbf{Основная идея.}
Обогатительная фабрика рассматривается как система потерь металла,
заданная матрицей
\[
  \text{металл} \times \text{класс крупности} \times \text{минеральная форма}
  \longrightarrow \text{потери, т}.
\]
Гипотеза улучшения не оценивается языковой моделью.
Она действует на конкретные ячейки матрицы, а эффект вычисляется
детерминированной моделью:
\[
  \text{эффект, т} =
  \text{потери в целевых ячейках}
  \times
  \text{коэффициент возврата},
\]
\[
  \text{эффект, \$} =
  \text{эффект, т}
  \times
  \text{цена соответствующего металла}.
\]
\end{mainbox}

\tableofcontents
\newpage

\section{Введение и постановка задачи}

Ключевое методологическое положение работы состоит в том, что коэффициент
возврата не является произвольным экспертным баллом. В базовой реализации он
хранится как параметр правила, однако математически представляет собой
дискретизацию физически обоснованной функции извлечения, зависящей от
крупности частиц, степени раскрытия минерала, типа потерь и доступного
оборудования, и уточняемую по результатам экспериментов.

\begin{formulaBox}
\[
\begin{aligned}
\text{правило модуля} &\equiv \text{формализованная инженерная гипотеза},\\
\text{коэффициент правила} &\equiv
  \text{калиброванная дискретизация кривой извлечения},\\
\text{ранг гипотезы} &\equiv
  \text{ожидаемый экономический эффект с учётом реализуемости,}\\
  &\quad \text{покрытия, тупиков и конфликтов.}
\end{aligned}
\]
\end{formulaBox}

\section{Обозначения}

Введём обозначения:
\begin{itemize}
  \item \(p \in P\) --- фабрика;
  \item \(m \in M\) --- металл, например \(\mathrm{Ni}\) или \(\mathrm{Cu}\);
  \item \(s \in S\) --- класс крупности, например \(+125\), \(-71{+}45\), \(-10\) (мкм);
  \item \(f \in F\) --- минеральная форма: свободная, запертая, рассеянная, нерентабельная;
  \item \(a \in A\) --- технологический модуль: доизмельчение, классификация, флотация тонких классов;
  \item \(c = (m,s,f)\) --- ячейка матрицы потерь; \(n = |M|\,|S|\,|F|\) --- число ячеек;
  \item \(x_{p,c} \ge 0\) --- тоннаж потерь металла в ячейке \(c\) на фабрике \(p\);
  \item \(q_m > 0\) --- цена металла \(m\) за тонну; \(m(c)\) --- металл ячейки \(c\);
  \item \(r_f \in \{0,1\}\) --- индикатор технологической извлекаемости формы \(f\);
  \item \(I_{p,a} \in \{0,1\}\) --- индикатор доступности оборудования модуля \(a\) на фабрике \(p\);
  \item \(k_{a,s,f} \in [0,1]\) --- коэффициент потенциального возврата модуля \(a\);
  \item \(h \in \Hyp\) --- гипотеза с множеством целевых ячеек \(C_h \subset M \times S \times F\).
\end{itemize}

Эффект гипотезы в тоннах и в денежном выражении:
\[
  T_h = \sum_{c \in C_h} x_c\, k_{h,c},
  \qquad
  U_h = \sum_{c \in C_h} x_c\, k_{h,c}\, q_{m(c)}.
\]
Индекс цены \(q_{m(c)}\) связан с ячейкой суммирования \(c\)
(в исходной записи \(q_m\) индекс \(m\) оставался несвязанным).

\section{Математический анализ: кривая извлечения}

Физической основой модели служит зависимость извлечения от крупности и
раскрытия: извлечение снижается на крупных классах из-за недостаточного
раскрытия минерала и на тонких классах из-за шламообразования и потери
селективности; максимум, как правило, достигается в среднем диапазоне
крупности.

Зависимость задаётся гладкой мультипликативной функцией
\[
  R_a(d,\lambda) =
  R_{\max,a} \cdot B_a(d) \cdot L_a(\lambda) \cdot E_{p,a},
\]
где:
\begin{itemize}
  \item \(d\) --- характерный размер частицы;
  \item \(\lambda \in [0,1]\) --- степень раскрытия полезного минерала;
  \item \(B_a(d) \in [0,1]\) --- безразмерная форма кривой по крупности;
  \item \(L_a(\lambda) \in [0,1]\) --- вклад раскрытия;
  \item \(E_{p,a} = I_{p,a}\,\varepsilon_{p,a}\) --- множитель реализуемости:
    индикатор наличия оборудования \(I_{p,a}\) и поправка на его состояние
    \(\varepsilon_{p,a} \in (0,1]\).
\end{itemize}

Для \(B_a(d)\) удобна колоколообразная кривая, гауссова в логарифмической
шкале крупности:
\[
  B_a(d) =
  \exp\!\left(
    -\frac{\bigl(\ln (d/d_{\mathrm{ref}}) - \mu_a\bigr)^2}{2\sigma_a^2}
  \right),
\]
где \(d_{\mathrm{ref}}\) --- опорный размер (например, 1~мкм), обеспечивающий
безразмерность аргумента логарифма. Отметим, что это не плотность
логнормального распределения (в ней присутствует множитель \(1/d\) и
нормировка), а безразмерная функция формы: она близка к нулю на слишком тонких
и слишком крупных классах и достигает максимума при
\(d = d_{\mathrm{ref}}\, e^{\mu_a}\) --- технологически оптимальной крупности
модуля \(a\).

Вклад раскрытия задаётся линейно:
\[
  L_a(\lambda) = \lambda,
\]
или с насыщением:
\[
  L_a(\lambda) = \frac{1 - \exp(-\beta_a \lambda)}{1 - \exp(-\beta_a)},
  \qquad \beta_a > 0.
\]
Нормировка во второй форме обеспечивает \(L_a(1) = 1\), так что
\(R_{\max,a}\) сохраняет смысл предельного извлечения.

Для дискретного класса крупности \(s = [d_{\min}, d_{\max}]\) коэффициент
правила получается усреднением по классу с учётом извлекаемости формы:
\[
  k_{a,s,f} =
  r_f \cdot
  \frac{
    \int_{d_{\min}}^{d_{\max}} R_a(d,\lambda_f)\, w_s(d)\dd d
  }{
    \int_{d_{\min}}^{d_{\max}} w_s(d)\dd d
  },
\]
где \(w_s(d) \ge 0\) --- плотность распределения частиц внутри класса,
\(\lambda_f\) --- характерная степень раскрытия формы \(f\). Множитель
\(r_f\) формализует обнуление коэффициента для технологически неизвлекаемой
формы.

\begin{mainbox}
В системе интегральное значение хранится как \code{rule.coeff}, а интервал
неопределённости --- как \code{rule.coeff\_min} и \code{rule.coeff\_max}.
Для каждого коэффициента прослеживается цепочка вывода:
\[
  \text{класс крупности}
  \to
  \text{форма минерала}
  \to
  \text{кривая извлечения}
  \to
  \text{коэффициент}
  \to
  \text{эффект}.
\]
\end{mainbox}

\section{Линейная алгебра: матрица потерь, покрытие и конфликты}

Матрица потерь фабрики разворачивается в вектор
\[
  x_p \in \R_+^n,
\]
где каждая координата соответствует одной ячейке
\(\text{металл} \times \text{класс} \times \text{форма}\).
Каждая гипотеза задаёт вектор действия
\[
  a_h \in [0,1]^n,
  \qquad
  a_{h,i} = k_{h,i}\ \text{при}\ i \in C_h,
  \qquad
  a_{h,i} = 0\ \text{иначе}.
\]

Координата \(a_{h,i}\) --- доля потерь ячейки \(i\), которую гипотеза способна
вернуть. Вектор эффекта и денежный эффект:
\[
  e_h = a_h \odot x_p,
  \qquad
  U_h = q^\top e_h,
\]
где \(\odot\) --- покоординатное произведение (произведение Адамара), а
\(q \in \R_+^n\) --- вектор цен, в котором ячейке \(i\) сопоставлена цена её
металла \(q_{m(i)}\).

Портфель гипотез \(S \subseteq \Hyp\) описывается матрицей эффектов
\[
  E = [\,e_1\ e_2\ \dots\ e_{|\Hyp|}\,] \in \R_+^{\,n \times |\Hyp|}
\]
(строки --- ячейки потерь, столбцы --- гипотезы). Покрытие ячейки \(i\):
\[
  \covered_i(S) =
  \min\Bigl(x_i,\ \sum_{h \in S} e_{h,i}\Bigr).
\]

Усечение минимумом принципиально: несколько гипотез не могут суммарно вернуть
из ячейки больше металла, чем в ней потеряно (запрет двойного счёта). Отсюда
формальное определение конфликта:
\[
  \conflict_i(S) =
  \mathbf{1}\Bigl[
    \sum_{h \in S} e_{h,i} > x_i
    \ \wedge\
    \bigl|\{h \in S : e_{h,i}>0\}\bigr| \ge 2
  \Bigr].
\]

В поисковом (RAG) слое линейная алгебра используется для поиска по корпусу:
\[
  \cos(u,v) =
  \frac{u^\top v}{\|u\|\,\|v\|}.
\]
Поисковый слой служит для подбора литературных обоснований и не участвует в
вычислении эффекта.

\section{Дифференциальные уравнения: кинетика флотации и динамика раскрытия}

Динамика извлечения описывается классической кинетической моделью флотации
первого порядка:
\[
  \frac{\mathrm{d}R}{\mathrm{d}t} =
  \kappa(d,\lambda)\bigl(R_\infty(d,\lambda)-R(t)\bigr),
  \qquad R(0) = 0.
\]

Решение при указанном начальном условии:
\[
  R(t) =
  R_\infty(d,\lambda)
  \bigl(1-\exp(-\kappa(d,\lambda)\,t)\bigr),
\]
где \(R(t)\) --- накопленное извлечение к моменту \(t\),
\(R_\infty \in [0,1]\) --- предельное извлечение для данного класса крупности
и степени раскрытия, \(\kappa > 0\) --- константа скорости флотации.
Модуль флотации тонких классов интерпретируется как изменение параметров
\(R_\infty\) и \(\kappa\) для тонких и рассеянных классов.

Для доизмельчения используется упрощённая динамика эффективной доли
раскрытого извлекаемого минерала:
\[
  \frac{\mathrm{d}\lambda}{\mathrm{d}\tau} =
  \gamma_a(d)(1-\lambda) - \rho_a(d)\lambda,
\]
где \(\tau\) --- условная интенсивность (доза) доизмельчения,
\(\gamma_a(d)\) --- скорость раскрытия запертых сростков,
\(\rho_a(d)\) --- скорость выбытия раскрытого минерала в неизвлекаемые шламы
при переизмельчении. Уравнение линейно и имеет устойчивое стационарное решение
\[
  \lambda^\ast = \frac{\gamma_a(d)}{\gamma_a(d)+\rho_a(d)},
\]
к которому \(\lambda(\tau)\) сходится экспоненциально со скоростью
\(\gamma_a+\rho_a\). Отсюда следует качественный вывод, важный для
маршрутизации гипотез: для крупных запертых сростков
\(\gamma_a \gg \rho_a\) и \(\lambda^\ast \to 1\) (доизмельчение эффективно),
тогда как для тонких классов \(\rho_a\) растёт и \(\lambda^\ast\) снижается
(доизмельчение контрпродуктивно). Модель сознательно не претендует на полноту
population-balance-подхода и используется как локальная инженерная
аппроксимация.

\section{Дифференциальная геометрия: пространство технологических состояний}

Дифференциально-геометрический язык применяется для формализации трёх
аспектов:
\begin{itemize}
  \item похожести фабрик;
  \item переносимости тупиков и правил;
  \item локальной чувствительности эффекта к изменению технологического состояния.
\end{itemize}

Фабрика представляется точкой на многообразии технологических состояний:
\[
  z_p = (\pi_p,\, X_p,\, e_p,\, \theta_p,\, b_p) \in \Manifold,
\]
где:
\begin{itemize}
  \item \(\pi_p \in \Delta^{n-1}\) --- нормированная матрица потерь (доли по ячейкам):
    \(\sum_i \pi_{p,i}=1\), \(\pi_{p,i}\ge 0\);
  \item \(X_p = \sum_i x_{p,i} > 0\) --- суммарный тоннаж потерь (масштаб),
    так что \(x_p = X_p\,\pi_p\);
  \item \(e_p\) --- вектор признаков оборудования;
  \item \(\theta_p\) --- текущие коэффициенты правил;
  \item \(b_p\) --- признаки лабораторной истории: успехи, частичные успехи, тупики.
\end{itemize}

Явное включение масштаба \(X_p\) устраняет несогласованность исходной
формулировки: целевая функция стоимости потерь определена на абсолютных
величинах, а не на долях. Пространство состояний --- произведение
\[
  \Manifold =
  \bigl(\Delta^{n-1} \times \R_+\bigr)
  \times
  E_{\text{equipment}}
  \times
  \Theta_{\text{rules}}
  \times
  B_{\text{history}}.
\]

\subsection{Метрика похожести фабрик}

Для сравнения фабрик вводится риманова метрика
\[
  g_z(u,v) = u^\top G(z)\,v,
  \qquad u,v \in T_z\Manifold,
\]
где \(T_z\Manifold\) --- касательное пространство в точке \(z\),
\(G(z)\) --- симметричная положительно определённая матрица весов признаков.
В базовой реализации \(G\) блочно-диагональна:
\[
  G =
  \diag
  (w_{\text{loss}},\, w_{\text{scale}},\, w_{\text{equipment}},\,
   w_{\text{rule}},\, w_{\text{history}}).
\]

Локальный элемент длины \(ds^2 = dz^\top G\, dz\); полное расстояние ---
геодезическая длина
\[
  d_\Manifold(p,p')
  =
  \inf_{\gamma}
  \int_0^1
  \sqrt{\gamma'(t)^\top G(\gamma(t))\,\gamma'(t)}
  \dd t.
\]

При постоянной \(G\) геодезические прямолинейны, и расстояние сводится к
практичной метрике Махаланобиса
\[
  d_G(p,p') =
  \sqrt{(z_p-z_{p'})^\top G\,(z_p-z_{p'})}.
\]

\emph{Оговорка.} Компонента \(\pi\) лежит на симплексе, поэтому евклидова
метрика на нём является лишь приближением; при необходимости более корректной
геометрии долей применимы метрика Фишера--Рао или геометрия Айчисона.
Для целей ранжирования и переноса опыта взвешенная евклидова аппроксимация
достаточна. Весовая матрица \(G\) кодирует экспертное знание: совпадение по
причине потерь, классу крупности и минеральной форме должно весить больше, чем
совпадение по суммарному тоннажу.

\subsection{Гипотеза как касательный вектор}

Гипотеза --- локальное действие на состояние фабрики, то есть вектор
касательного пространства:
\[
  V_h(z_p) \in T_{z_p}\Manifold,
\]
компоненты которого в координатах потерь равны \(-e_{h,i} = -x_i k_{h,i}\)
для \(i \in C_h\). Если целевая функция --- стоимость потерь
\[
  J(z) = \sum_i q_i x_i = X_p \sum_i q_i \pi_{p,i},
\]
то мгновенная польза гипотезы есть взятая с обратным знаком производная
\(J\) по направлению \(V_h\):
\[
  \operatorname{Benefit}(h)
  =
  -\,dJ_z(V_h)
  =
  -\nabla J(z_p)^\top V_h(z_p)
  =
  \sum_{i \in C_h} q_i x_i k_{h,i}
  = U_h.
\]

Последнее равенство показывает, что дискретный расчёт движка в точности
реализует эту направленную производную: линейность \(J\) по \(x\) делает
формулу точной, а не приближённой.

\begin{mainbox}
Гипотеза --- не текстовая идея, а вектор улучшения на многообразии
технологических состояний фабрики.
\end{mainbox}

\subsection{Перенос тупиков между фабриками}

Если гипотеза \(h\) не подтвердилась на фабрике \(p\), предупреждение
переносится на фабрику \(p'\) при выполнении двух условий:
\[
  d_\Manifold(p,p') < \varepsilon
\]
и совпадении локальной области действия:
\[
  \operatorname{module}(h),\quad
  \operatorname{target\_cause}(h),\quad
  \operatorname{size\_class}(h).
\]

Тупик переносится не глобально, а локально: между близкими точками
многообразия и вдоль близких направлений. В базовой реализации критерий
упрощён до совпадения дискретных атрибутов; математическим содержанием
остаётся принцип локальной переносимости опыта между близкими
технологическими состояниями.

\subsection{Чувствительность и второй порядок}

Локальная чувствительность оценки гипотезы к возмущениям состояния (ошибки
измерения потерь, нестабильность питания):
\[
  S_h(z) = \|\nabla_z \Score_h(z)\|_{G^{-1}},
  \qquad
  \|v\|_{G^{-1}} = \sqrt{v^\top G^{-1} v}.
\]
Двойственная норма \(G^{-1}\) для градиента согласована с нормой \(G\) на
приращениях состояния:
\(|\Score_h(z+u)-\Score_h(z)| \le
\|\nabla\Score_h\|_{G^{-1}}\,\|u\|_G + o(\|u\|)\).
Большое значение \(S_h(z)\) означает высокий эффект при высокой
чувствительности к данным, что интерпретируется как риск.

Второй порядок описывается гессианом
\(H_h(z) = \nabla^2 \Score_h(z)\): при большой его норме линейная
экстраполяция оценки между фабриками ненадёжна, что даёт формальное
обоснование осторожного переноса правил. Явное вычисление гессиана в базовой
реализации не требуется.

\subsection{Геометрия корпуса документов}

Нормированные эмбеддинги документов лежат на единичной сфере:
\[
  \frac{v}{\|v\|} \in S^{d-1}.
\]

Косинусная близость монотонно связана с геодезическим (угловым) расстоянием
на сфере:
\[
  d_S(u,v) = \arccos(u^\top v),
  \qquad \|u\| = \|v\| = 1.
\]

Эта геометрия применяется исключительно для подбора обоснований и не
подменяет численный расчёт эффекта.

\section{Теория вероятностей: неопределённость эффекта и риск}

Коэффициент возврата известен неточно и моделируется случайной величиной
\(K_{a,s,f}\) с носителем на экспертном интервале
\([k_{\min}, k_{\max}] \subseteq [0,1]\). Хранимая пара
\(k_{\min} \le k \le k_{\max}\) трактуется как границы носителя (или
доверительного интервала) априорного распределения. Удобная параметризация ---
масштабированное бета-распределение:
\[
  K = k_{\min} + (k_{\max}-k_{\min})\,Z,
  \qquad Z \sim \operatorname{Beta}(\alpha,\beta),
\]
с моментами
\[
  \E[K] = k_{\min} + (k_{\max}-k_{\min})\frac{\alpha}{\alpha+\beta},
  \qquad
  \operatorname{Var}[K] =
  (k_{\max}-k_{\min})^2
  \frac{\alpha\beta}{(\alpha+\beta)^2(\alpha+\beta+1)}.
\]

При независимости коэффициентов по ячейкам ожидаемые эффекты линейны:
\[
  \E[T_h] =
  \sum_{c \in C_h} x_c\, \E[K_{h,c}],
  \qquad
  \E[U_h] =
  \sum_{c \in C_h} x_c\, q_{m(c)}\, \E[K_{h,c}].
\]

Помимо математического ожидания, вычисляется вероятность превышения порога:
\[
  \Prob(U_h > U_{\min}).
\]

Это существенно для планирования лабораторного этапа: гипотеза с большим
средним эффектом, но высокой дисперсией может уступить более скромной, но
надёжной. Вероятностная модель задаёт семантику статусов эксперимента:
\begin{itemize}
  \item \code{success} --- наблюдение попало в ожидаемый диапазон или выше;
  \item \code{partial} --- эффект обнаружен, но ниже ожидаемого;
  \item \code{failure} --- эффект статистически не подтверждён;
    формируется кандидат в тупик.
\end{itemize}

\section{Математическая статистика: калибровка коэффициентов}

После лабораторного эксперимента система получает наблюдение
\[
  y_n =
  \clip\!\left(
  \frac{\text{measured effect}}{\text{predicted max effect}},\,0,\,1\right),
\]
либо непосредственно измеренный коэффициент возврата \(y_n \in [0,1]\).
Усечение к \([0,1]\) обеспечивает корректность области значений.
Калибровка выполняется рекуррентным взвешенным усреднением:
\[
  k_{n+1}
  =
  \frac{N k_n + y_n}{N+1},
\]
где \(k_n\) --- текущий коэффициент, \(N\) --- число накопленных калибровок.
Эквивалентная форма стохастической аппроксимации:
\[
  k_{n+1}
  =
  k_n + \eta_n(y_n-k_n),
  \qquad
  \eta_n = \frac{1}{N+1}.
\]

Шаг \(\eta_n = 1/(N+1)\) соответствует несмещённому выборочному среднему;
более консервативное обновление:
\[
  \eta_n = \frac{\eta_0}{\sqrt{N+1}}.
\]
Оба режима удовлетворяют условиям Роббинса--Монро
(\(\sum_n \eta_n = \infty\), \(\sum_n \eta_n^2 < \infty\)), гарантирующим
сходимость к среднему значению наблюдений при стационарном процессе.
Обновление допускает и байесовскую интерпретацию: при бета-априорном
распределении рекуррентная формула воспроизводит эволюцию апостериорного
среднего.

После обновления коэффициент проецируется на экспертный интервал:
\[
  k_{n+1}
  \leftarrow
  \clip(k_{n+1}, k_{\min}, k_{\max}).
\]

Для эксперимента со статусом \code{failure} применяется консервативная
политика
\[
  k_{n+1} \leftarrow k_{\min}
\]
с одновременным созданием записи в базе тупиков:
\[
  \operatorname{dead\_end}
  (\operatorname{module},\,\operatorname{target\_cause},\,
   \operatorname{size\_class},\,\operatorname{reason}).
\]

Каждый эксперимент либо уточняет коэффициент, либо фиксирует воспроизводимый
отрицательный результат, переносимый на близкие фабрики. Обучение остаётся
полностью интерпретируемым.

\section{Ранжирование гипотез}

Для гипотезы \(h\) определены величины:
\begin{itemize}
  \item \(T_h\) --- эффект в тоннах;
  \item \(U_h\) --- эффект в денежном выражении;
  \item \(F_h\) --- индикатор реализуемости по оборудованию;
  \item \(D_h\) --- индикатор тупика;
  \item \(\Gamma_h\) --- индикатор конфликта за ячейки.
\end{itemize}
Обозначение конфликта заменено с \(C_h\) на \(\Gamma_h\) для устранения
коллизии с множеством целевых ячеек \(C_h\).

Базовая реализация использует прозрачный аддитивный скоринг:
\[
  \operatorname{score}_h
  =
  \min\!\left(\frac{U_h}{100\,000},\,70\right)
  +
  \min(T_h,\,20)
  +
  10\,\mathbf{1}[F_h]
\]
с последующей лексикографической сортировкой:
\[
  \text{реализуемость}
  \to
  \operatorname{score}
  \to
  \text{денежный эффект}.
\]
Слагаемые безразмерны за счёт фиксированных нормировочных констант;
насыщение через минимум ограничивает вклад каждого критерия.

Расширенная версия скоринга:
\[
\begin{aligned}
  \Score_h
  ={}&
  w_u \norm(\E[U_h])
  + w_t \norm(\E[T_h])
  + w_p \Prob(U_h > U_{\min})\\
  &+ w_{\text{cov}}\,\norm(\Delta\Cov_h)
  - w_r \Risk_h
  - w_d \mathbf{1}[D_h]
  - w_c \operatorname{ConflictPenalty}_h,
\end{aligned}
\]
где \(\norm(\cdot)\) --- нормировка критериев к сопоставимой шкале
(например, min--max по текущему пулу гипотез), а маргинальный вклад в
покрытие при текущем портфеле \(S\):
\[
  \Delta\Cov_h(S)
  =
  \sum_i \min\bigl(x_i,\,\covered_i(S)+e_{h,i}\bigr)
  -
  \sum_i \covered_i(S).
\]

Слагаемое покрытия поощряет гипотезы, закрывающие ещё не охваченные потери,
и предотвращает конкуренцию нескольких гипотез за одну ячейку при
игнорировании соседних крупных потерь.

\section{Оптимизация портфеля экспериментов}

Выбор портфеля гипотез формулируется как задача комбинаторной оптимизации
\[
  \max_{S \subseteq \Hyp} F(S),
  \qquad
  F(S) = g(S) - \sum_{h \in S} \operatorname{cost}_h
  - \operatorname{penalties}(S),
\]
где ценность покрытия
\[
  g(S)
  =
  \sum_i q_i
  \min\Bigl(x_i,\,\sum_{h \in S} e_{h,i}\Bigr)
\]
при ограничениях
\[
  \sum_{h \in S} \operatorname{cost}_h \le \operatorname{Budget},
  \qquad
  \sum_{h \in S} \operatorname{duration}_h \le \operatorname{TimeLimit},
  \qquad
  \text{ограничения по оборудованию}.
\]

\begin{mainbox}
\textbf{Утверждение (субмодулярность покрытия).}
Функция \(g\) монотонно неубывающая и субмодулярная: для любых
\(S \subseteq S' \subseteq \Hyp\) и \(h \notin S'\)
\[
  g(S \cup \{h\}) - g(S)
  \ \ge\
  g(S' \cup \{h\}) - g(S').
\]
\emph{Доказательство} сводится к покоординатной проверке: при фиксированном
\(i\) функция \(t \mapsto \min(x_i, t)\) вогнута и неубывающая, а
суперпозиция вогнутой неубывающей функции с аддитивной функцией множества
субмодулярна; неотрицательная линейная комбинация субмодулярных функций
субмодулярна.
\end{mainbox}

Субмодулярность формализует эффект убывающей отдачи (почти закрытая ячейка
даёт малый прирост следующей гипотезе) и обосновывает жадный алгоритм:
\[
  h^\ast
  =
  \arg\max_{h \notin S}
  \frac{\Delta F(h\mid S)}{\operatorname{cost}_h}.
\]

Для монотонных субмодулярных функций при бюджетном ограничении жадный
алгоритм по отношению <<прирост/стоимость>> (в комбинации с выбором лучшего
одиночного элемента) даёт гарантию приближения
\(1 - 1/e \approx 0{,}632\) от оптимума; практичная эвристика получает
доказуемую нижнюю оценку качества.

Если дорожные карты гипотез содержат совпадающие лабораторные этапы,
стоимость портфеля агрегируется по уникальным этапам:
\[
  \operatorname{Cost}(S)
  =
  \sum_{\operatorname{unique}\ \operatorname{shared\_key}}
  \operatorname{cost}_{\operatorname{step}},
\]
что превращает результат из простого рейтинга в план экспериментов с
объединением общих лабораторных действий. Агрегирование стоимости делает
\(F\) в общем случае не субмодулярной; гарантия \(1-1/e\) относится к
покрывающей части \(g\), а объединение этапов лишь улучшает достижимое
значение.

\section{Роль разделов математики в модели}

\begin{longtable}{p{0.28\linewidth}p{0.64\linewidth}}
\toprule
\textbf{Раздел} & \textbf{Роль в модели}\\
\midrule
\endhead

Математический анализ &
Гладкая функция извлечения от крупности и раскрытия; обоснование
маршрутизации: крупные запертые сростки --- в доизмельчение, тонкие
рассеянные классы --- в тонкую флотацию.\\
\addlinespace

Линейная алгебра &
Векторизация матрицы потерь, векторы действия гипотез, денежный эффект,
покрытие и формальный критерий конфликта.\\
\addlinespace

Дифференциальные уравнения &
Кинетика флотации первого порядка и динамика раскрытия при доизмельчении
со стационарным анализом.\\
\addlinespace

Дифференциальная геометрия &
Фабрики как точки многообразия состояний, гипотезы как касательные векторы,
метрика похожести, локальный перенос тупиков, чувствительность и геометрия
эмбеддингов.\\
\addlinespace

Теория вероятностей &
Неопределённость коэффициентов (масштабированное бета-распределение),
ожидаемый эффект, вероятность превышения порога, риск.\\
\addlinespace

Математическая статистика &
Рекуррентная калибровка коэффициентов по экспериментам с гарантиями
сходимости и байесовской интерпретацией.\\
\bottomrule
\end{longtable}

\section{Заключение}

Построена математическая модель диагностики потерь металлов и ранжирования
гипотез, объединяющая шесть разделов математики в единую вычислительную
схему. Основные результаты состоят в следующем.

\begin{enumerate}
  \item \textbf{Детерминированность и прослеживаемость расчёта.}
    Эффект каждой гипотезы вычисляется по явным формулам из ячеек матрицы
    потерь и коэффициентов правил; языковая модель не участвует в численной
    оценке и используется лишь для работы с текстом и литературой.
    Для каждого коэффициента прослеживается цепочка вывода от класса
    крупности и минеральной формы до денежного эффекта.
  \item \textbf{Физическая обоснованность коэффициентов.}
    Коэффициент возврата определён как усреднение по классу крупности
    физически осмысленной кривой извлечения, зависящей от крупности,
    раскрытия и оборудования; кинетические уравнения флотации и раскрытия
    объясняют направление действия модулей, включая стационарный анализ,
    показывающий границы применимости доизмельчения.
  \item \textbf{Корректность агрегирования.}
    Усечение покрытия минимумом исключает двойной счёт потерь; конфликт
    гипотез за ячейку определён формально. Доказанная субмодулярность
    функции покрытия даёт жадному алгоритму формирования портфеля гарантию
    приближения \(1-1/e\), переводя практическую эвристику в разряд
    алгоритмов с доказуемым качеством.
  \item \textbf{Управление неопределённостью и обучение.}
    Коэффициенты трактуются как случайные величины с масштабированным
    бета-распределением на экспертном интервале; ранжирование учитывает
    не только ожидаемый эффект, но и вероятность превышения порога и
    чувствительность к данным. Рекуррентная калибровка удовлетворяет
    условиям Роббинса--Монро и допускает байесовскую интерпретацию,
    а отрицательные результаты фиксируются как переносимые тупики.
  \item \textbf{Масштабируемость через геометрию состояний.}
    Представление фабрик точками многообразия технологических состояний
    с взвешенной метрикой формализует похожесть фабрик и локальный перенос
    опыта: правила и тупики переносятся только между близкими состояниями
    и близкими направлениями действия.
\end{enumerate}

Научная ценность модели заключается не в построении полномасштабной
first-principles модели обогащения, а в методически выверенном компромиссе:
известная форма физической зависимости, калибровка на данных конкретной
фабрики, прозрачный матричный расчёт эффекта, статистическое обучение по
экспериментам и структурный запрет двойного счёта потерь. Такая конструкция
достаточно строга для экспертизы техническим комитетом и достаточно проста
для полной верификации всех вычислений вручную.

Перспективными направлениями развития являются: переход от евклидовой
аппроксимации к геометрии Айчисона или метрике Фишера--Рао на симплексе
долей потерь; замена точечной калибровки полноценным байесовским обновлением
параметров бета-распределения; учёт корреляций коэффициентов между ячейками;
постановка выбора портфеля как задачи субмодулярной максимизации при
матроидных ограничениях с соответствующими гарантиями.

\end{document}
